% modul-1.tex - Modul 1: Programmier-Mindset & KI-Tools Setup

\startcomponent modul-1

\chapter{Programmier-Mindset \& KI-Tools Setup}

{\tfx\it Dauer: 1 Tag (4 Lektionen à 50 Minuten) · Voraussetzung: Keine}

\blank[big]

\startlernziele
{\bf Lernziele}

Nach diesem Modul können Sie:
\startitemize[n,packed]
    \item Erklären, was Programmieren im Kern bedeutet und wie sich KI-gestütztes Programmieren unterscheidet
    \item Eine vollständige Entwicklungsumgebung (Python, Git, VS Code, uv) eigenständig installieren
    \item KI-Assistenten (GitHub Copilot, ChatGPT, Claude) effektiv für einfache Programmieraufgaben nutzen
    \item Die Qualität von KI-generiertem Code auf grundlegende Fehler überprüfen
    \item Mit Hilfe von Prompts einfache Python-Programme generieren lassen
    \item Beurteilen, wann der Einsatz von KI-Tools sinnvoll ist
    \item Den typischen Workflow moderner Softwareentwicklung mit Git beschreiben
    \item Effektive Prompts für KI-Coding-Assistenten formulieren und iterativ verbessern
\stopitemize
\stoplernziele

% =============================================================================
% VORBEREITUNG
% =============================================================================

\section{Vorbereitung}

{\it Zeitaufwand: 2--3 Stunden · Deadline: Vor der ersten Präsenzveranstaltung}

\blank[medium]

Bevor wir mit dem Präsenzunterricht starten, bereiten Sie bitte Ihre Entwicklungsumgebung vor und machen sich mit den Grundkonzepten vertraut.

\subsection{Entwicklungsumgebung wählen}

Wählen Sie {\bf EINE} der folgenden Optionen:

\subsubsection{Option A: GitHub Codespaces (empfohlen)}

\startitemize[packed]
    \item Keine Installation nötig -- läuft im Browser
    \item GitHub Student Developer Pack aktivieren
    \item Zeitaufwand: ca. 15 Minuten
    \item Voraussetzung: .edu E-Mail-Adresse
    \item Kosten: Kostenlos (60 Stunden/Monat)
\stopitemize

\subsubsection{Option B: Lokale Installation}

\startitemize[packed]
    \item Python installiert und getestet
    \item Git installiert und konfiguriert
    \item VS Code installiert und eingerichtet
    \item uv installiert
    \item GitHub Copilot oder Alternative eingerichtet
    \item Zeitaufwand: ca. 90 Minuten
    \item Voraussetzung: Admin-Rechte auf Ihrem Computer
\stopitemize

\startinfobox
{\bf Tipp: Hybrid-Ansatz}

Nutzen Sie beide Optionen! Codespaces für Kursmaterialien und Übungen, lokal für grössere Projekte und Experimente. So sparen Sie Codespaces-Stunden und haben maximale Flexibilität.
\stopinfobox

\subsection{Installationsanleitung}

{\it Zeitaufwand: ca. 90 Minuten}

\subsubsection{1. Python installieren}

{\bf Windows:}
\startitemize[packed]
    \item Besuchen Sie python.org/downloads
    \item Laden Sie Python 3.11 oder höher herunter
    \item {\bf WICHTIG:} Aktivieren Sie "Add Python to PATH" während der Installation
\stopitemize

{\bf macOS (mit Homebrew):}
\startbashcode
# Homebrew installieren (falls noch nicht vorhanden)
/bin/bash -c "$(curl -fsSL https://raw.githubusercontent.com/Homebrew/install/HEAD/install.sh)"

# Python installieren
brew install python@3.11
\stopbashcode

{\bf Linux (Ubuntu/Debian):}
\startbashcode
sudo apt update
sudo apt install python3.11 python3.11-venv python3-pip
\stopbashcode

{\bf Testen:}
\startbashcode
python --version   # oder python3 --version
# Sollte ausgeben: Python 3.11.x oder höher
\stopbashcode

\subsubsection{2. Git installieren}

{\bf Windows:} Laden Sie den Installer von git-scm.com herunter.

{\bf macOS:}
\startbashcode
brew install git
# ODER: Xcode Command Line Tools
xcode-select --install
\stopbashcode

{\bf Linux:}
\startbashcode
sudo apt install git
\stopbashcode

{\bf Git konfigurieren (alle Systeme):}
\startbashcode
git config --global user.name "Ihr Name"
git config --global user.email "ihre.email@example.com"
\stopbashcode

\subsubsection{3. VS Code installieren}

\startitemize[n,packed]
    \item Besuchen Sie code.visualstudio.com
    \item Laden Sie die Version für Ihr Betriebssystem herunter
    \item Installieren Sie VS Code
\stopitemize

{\bf Empfohlene Extensions:}
\startitemize[packed]
    \item {\bf Python} (Microsoft) -- Extension ID: \code{ms-python.python}
    \item {\bf GitHub Copilot} -- Extension ID: \code{GitHub.copilot}
    \item {\bf GitLens} (optional) -- Extension ID: \code{eamodio.gitlens}
\stopitemize

\subsubsection{4. uv installieren}

{\bf uv} ist ein modernes, schnelles Tool für Python-Paketverwaltung.

{\bf Windows (PowerShell):}
\startbashcode
powershell -c "irm https://astral.sh/uv/install.ps1 | iex"
\stopbashcode

{\bf macOS/Linux:}
\startbashcode
curl -LsSf https://astral.sh/uv/install.sh | sh
\stopbashcode

\subsubsection{5. KI-Assistent einrichten}

{\bf Option A: GitHub Copilot (empfohlen)}
\startitemize[packed]
    \item GitHub-Account erstellen (falls noch nicht vorhanden)
    \item Studenten: Kostenloser Zugang über education.github.com/pack
    \item In VS Code anmelden (GitHub-Symbol unten links)
\stopitemize

{\bf Option B: Alternative KI-Tools}
\startitemize[packed]
    \item {\bf ChatGPT:} Kostenloser Account auf chat.openai.com
    \item {\bf Claude:} Kostenloser Account auf claude.ai
\stopitemize

\subsection{Erste Schritte mit KI}

{\it Zeitaufwand: 30 Minuten}

Nutzen Sie ChatGPT oder Claude, um Ihre ersten Python-Programme zu erstellen.

\subsubsection{Schritt 1: "Hello World" generieren}

{\bf Ihr Prompt:}
\startbashcode
Erstelle ein einfaches "Hello World" Programm in Python.
\stopbashcode

Kopieren Sie den generierten Code und speichern Sie ihn als \code{hello.py}. Führen Sie ihn aus mit: \code{python hello.py}

\subsubsection{Schritt 2: Variante mit Benutzereingabe}

{\bf Ihr Prompt:}
\startbashcode
Erstelle eine Variante, die den Benutzer nach seinem Namen fragt
und dann "Hello, [Name]!" ausgibt.
\stopbashcode

{\bf Erwartetes Ergebnis:}
\startpythoncode
name = input("Wie ist dein Name? ")
print(f"Hello, {name}!")
\stoppythoncode

\subsubsection{Reflexion}

Beantworten Sie folgende Fragen:
\startitemize[packed]
    \item Wie genau waren die KI-Antworten? Hat der Code auf Anhieb funktioniert?
    \item Wie hilfreich waren die Erklärungen?
    \item Wie würden Sie die Prompts verbessern?
\stopitemize

\subsection{Leseauftrag}

{\it Zeitaufwand: 45 Minuten}

\subsubsection{Materialien}

\startitemize[n,packed]
    \item {\bf Artikel:} "What is Programming?" (15 Min.) -- freeCodeCamp oder KI-Erklärung
    \item {\bf Video:} "How AI is Changing Programming" (15 Min.) -- YouTube-Suche
    \item {\bf Blog-Post:} "Effective Prompting for Code Generation" (15 Min.) -- OpenAI Prompt Engineering Guide
\stopitemize

\subsubsection{Reflexionsfragen}

\startitemize[n,packed]
    \item Was überrascht Sie am meisten an der modernen Programmierung?
    \item Welche Bedenken haben Sie bezüglich KI-gestütztem Programmieren?
    \item Was erwarten Sie von diesem Kurs zu lernen?
    \item Warum möchten Sie programmieren lernen?
\stopitemize

% =============================================================================
% PRAXIS
% =============================================================================

\section{Praxis}

{\it Präsenzunterricht: 4 Lektionen à 50 Minuten}

\subsection{Lektion 1: Einführung und Mindset}

{\it Dauer: 50 Minuten}

\startinfobox
{\bf Kernbotschaft}

Programmieren ist Problemlösung, nicht Syntax-Memorierung!
\stopinfobox

\subsubsection{Die 4 Säulen des Programmierens}

{\bf 1. Problemdekomposition}

Grosse Probleme in kleinere, lösbare Teilprobleme zerlegen.

{\it Beispiel: "Kaffee kochen"}
\startitemize[n,packed]
    \item Wasser holen
    \item Wasser in Maschine füllen
    \item Kaffeepulver einfüllen
    \item Maschine einschalten
    \item Warten
    \item Kaffee einschenken
\stopitemize

{\bf 2. Algorithmisches Denken}

Schritt-für-Schritt-Anweisungen erstellen -- klare, eindeutige Abfolge, wiederholbare Prozesse.

{\bf 3. Abstraktion}

Unwichtige Details weglassen, auf das Wesentliche fokussieren.

{\it Beispiel: "Auto fahren"} -- Sie müssen nicht wissen, wie der Motor funktioniert. Sie müssen nur wissen: Gas, Bremse, Lenkrad.

{\bf 4. Musterkennung}

Ähnliche Probleme erkennen, bekannte Lösungen wiederverwenden.

\subsubsection{Die Rolle von KI}

{\bf KI als Werkzeug:}
\startitemize[packed]
    \item Generiert Code basierend auf Beschreibungen
    \item Erklärt komplexen Code
    \item Schlägt Verbesserungen vor
    \item Hilft beim Debugging
\stopitemize

{\bf Was KI NICHT ersetzt:}
\startitemize[packed]
    \item Problemverständnis
    \item Kreative Lösungsansätze
    \item Entscheidungen treffen
    \item Code-Qualität bewerten
\stopitemize

\subsection{Lektion 2: Setup-Vertiefung und Git Basics}

{\it Dauer: 50 Minuten}

\subsubsection{VS Code Grundlagen}

{\bf Die wichtigsten Bereiche:}
\startitemize[packed]
    \item {\bf Explorer} (links): Dateien und Ordner
    \item {\bf Editor} (Mitte): Code schreiben
    \item {\bf Terminal} (unten): Befehle ausführen
    \item {\bf Extensions} (links): Erweiterungen installieren
\stopitemize

{\bf Wichtige Shortcuts:}

\starttabulate[|l|l|l|]
\HL
\NC {\bf Aktion} \NC {\bf Windows/Linux} \NC {\bf macOS} \NC\NR
\HL
\NC Command Palette \NC Ctrl+Shift+P \NC Cmd+Shift+P \NC\NR
\NC Terminal öffnen \NC Ctrl+ö \NC Ctrl+ö \NC\NR
\NC Datei suchen \NC Ctrl+P \NC Cmd+P \NC\NR
\NC Speichern \NC Ctrl+S \NC Cmd+S \NC\NR
\HL
\stoptabulate

\subsubsection{Git-Grundlagen}

{\bf Was ist Git?}

Versionskontrolle -- speichert Änderungen an Dateien über Zeit, ermöglicht Zusammenarbeit, macht Änderungen rückgängig.

{\bf Grundlegende Git-Befehle:}

\startbashcode
git init                    # Neues Repository erstellen
git status                  # Status prüfen
git add dateiname.py        # Datei zum Staging hinzufügen
git add .                   # Alle Dateien hinzufügen
git commit -m "Nachricht"   # Commit erstellen
git log --oneline           # Historie anzeigen
\stopbashcode

{\bf Der typische Workflow:}
\startbashcode
# 1. Dateien ändern
# 2. Status prüfen
git status
# 3. Dateien zum Staging hinzufügen
git add .
# 4. Commit erstellen
git commit -m "feat: Neue Funktion hinzugefügt"
# 5. Wiederholen
\stopbashcode

\subsection{Lektion 3: Effektives Prompting}

{\it Dauer: 50 Minuten}

\startinfobox
{\bf Kernbotschaft}

Ein guter Prompt ist der Schlüssel zu gutem Code!
\stopinfobox

\subsubsection{Die 5 Elemente eines guten Prompts}

{\bf 1. Klarheit} -- Eindeutige, präzise Formulierung
\startitemize[packed]
    \item Schlecht: "Sortiere die Liste"
    \item Gut: "Sortiere die Liste von Zahlen aufsteigend"
\stopitemize

{\bf 2. Kontext} -- Hintergrundinformationen bereitstellen
\startitemize[packed]
    \item Schlecht: "Erstelle eine Funktion zum Validieren"
    \item Gut: "Erstelle eine Python-Funktion, die eine E-Mail-Adresse validiert..."
\stopitemize

{\bf 3. Spezifität} -- Genaue Anforderungen, erwartete Ein-/Ausgaben

{\bf 4. Beispiele} -- Konkrete Beispiele für Ein- und Ausgaben

{\bf 5. Einschränkungen} -- Was NICHT gemacht werden soll

\subsubsection{Prompt-Template}

\startbashcode
[Aufgabe]: Was soll gemacht werden?
[Kontext]: Warum? Wofür?
[Anforderungen]: Spezifische Details
[Beispiel]: Ein- und Ausgabe
[Einschränkungen]: Was zu beachten ist
\stopbashcode

\subsubsection{Häufige Fehler}

\startitemize[packed]
    \item {\bf Zu vage:} "Schreib mir was mit Daten"
    \item {\bf Zu komplex:} Alles auf einmal wollen -- zerlegen Sie grosse Aufgaben!
    \item {\bf Fehlender Kontext:} "Wie sortiere ich das?" (Was? Welche Sprache?)
    \item {\bf Keine Beispiele:} "Formatiere den String" (Wie genau?)
\stopitemize

\subsection{Lektion 4: Von der Idee zur App}

{\it Dauer: 50 Minuten}

\subsubsection{Der Entwicklungsworkflow}

\startitemize[n,packed]
    \item Problem verstehen
    \item In Teilprobleme zerlegen
    \item Für jedes Teilproblem: Prompt → Code → Verstehen → Testen → Anpassen
    \item Teillösungen zusammenführen
    \item Gesamtes Programm testen
    \item Code committen (Git)
\stopitemize

\subsubsection{Code-Review: Was prüfen?}

{\bf 1. Funktionalität} -- Macht der Code, was er soll?

{\bf 2. Lesbarkeit} -- Sind Variablennamen aussagekräftig?
\startpythoncode
# Schlecht
def f(x, y):
    return x * y * 0.19

# Gut
def calculate_tax(price, quantity):
    tax_rate = 0.19
    return price * quantity * tax_rate
\stoppythoncode

{\bf 3. Fehlerbehandlung} -- Werden Fehler abgefangen?
\startpythoncode
# Ohne Fehlerbehandlung
age = int(input("Alter: "))  # Crash bei "abc"

# Mit Fehlerbehandlung
try:
    age = int(input("Alter: "))
except ValueError:
    print("Bitte eine Zahl eingeben!")
    age = 0
\stoppythoncode

\subsubsection{Debugging-Strategien}

\startitemize[n,packed]
    \item {\bf Print-Debugging:} Werte an kritischen Stellen ausgeben
    \item {\bf Fehler lesen:} Fehlermeldungen verstehen lernen
    \item {\bf Schrittweise testen:} Kleine Teile einzeln prüfen
    \item {\bf KI um Hilfe fragen:} Code + Fehlermeldung an KI senden
\stopitemize

% =============================================================================
% ÜBUNGEN
% =============================================================================

\section{Übungen}

\startaufgabenbox
{\bf Übung 1: Einkaufsliste durchdenken} {\it (15 Min.)}

Zerlegen Sie das Problem "Einkaufsliste verwalten" in Teilprobleme:
\startitemize[packed]
    \item Welche Daten müssen gespeichert werden?
    \item Wie funktioniert "Artikel hinzufügen"?
    \item Was passiert beim "Artikel entfernen", wenn der Artikel nicht existiert?
    \item In welcher Reihenfolge werden Artikel angezeigt?
\stopitemize
\stopaufgabenbox

\blank[medium]

\startaufgabenbox
{\bf Übung 2: Git-Repository erstellen} {\it (20 Min.)}

\startitemize[n,packed]
    \item Projekt-Ordner erstellen: \code{mkdir mein-python-projekt}
    \item Git initialisieren: \code{git init}
    \item README.md mit KI erstellen
    \item .gitignore für Python erstellen
    \item Ersten Commit machen: \code{git commit -m "docs: Initiales Setup"}
\stopitemize
\stopaufgabenbox

\blank[medium]

\startaufgabenbox
{\bf Übung 3: Prompting-Challenge} {\it (20 Min.)}

Formulieren Sie vollständige Prompts für:
\startitemize[n,packed]
    \item {\bf Einkaufsliste:} Artikel hinzufügen, entfernen, anzeigen
    \item {\bf Temperatur-Umrechner:} Celsius ↔ Fahrenheit
    \item {\bf Wort-Zähler:} Wörter, Zeichen, längstes Wort
\stopitemize
\stopaufgabenbox

\blank[medium]

\startaufgabenbox
{\bf Übung 4: Einfacher Taschenrechner} {\it (25 Min.)}

Erstellen Sie einen Taschenrechner mit:
\startitemize[packed]
    \item Grundrechenarten (+, -, *, /)
    \item Benutzereingabe für zwei Zahlen
    \item Division durch Null abfangen
    \item Beenden mit 'quit'
\stopitemize

{\bf Testfälle:} 5+3=8, 10/0→Fehler, "abc"→Fehler
\stopaufgabenbox

% =============================================================================
% NACHBEARBEITUNG
% =============================================================================

\section{Nachbearbeitung}

{\it Hausaufgaben vor Modul 2}

\subsection{Aufgabe 1: Prompt-Portfolio}

Sammeln Sie 10 effektive Prompts:
\startitemize[packed]
    \item Dokumentieren Sie, was funktioniert hat
    \item Notieren Sie Verbesserungen
    \item Kategorisieren Sie nach Aufgabentyp
\stopitemize

\subsection{Aufgabe 2: Persönliches Projekt-Setup}

Erstellen Sie ein Git-Repository für Ihr Lernprojekt:
\startitemize[packed]
    \item README.md mit Projektbeschreibung
    \item .gitignore für Python
    \item Erste Python-Datei mit "Hello World"
    \item Mindestens 3 Commits
\stopitemize

\subsection{Aufgabe 3: Code-Review-Übung}

Überprüfen Sie KI-generierten Code auf:
\startitemize[packed]
    \item Funktionalität (funktioniert er?)
    \item Lesbarkeit (verständlich?)
    \item Fehlerbehandlung (was kann schiefgehen?)
\stopitemize

\subsection{Reflexion}

Beantworten Sie schriftlich:
\startitemize[n,packed]
    \item Was war die wichtigste Erkenntnis aus Modul 1?
    \item Was hat gut funktioniert beim Arbeiten mit KI?
    \item Welche Herausforderungen sind aufgetreten?
    \item Was möchten Sie in Modul 2 vertiefen?
\stopitemize

% =============================================================================
% MATERIALIEN
% =============================================================================

\section{Materialien}

\subsection{Handout: Programmier-Mindset}

Die 4 Säulen des Programmierens:
\startitemize[n,packed]
    \item {\bf Problemdekomposition} -- Grosse Probleme zerlegen
    \item {\bf Algorithmisches Denken} -- Schritt-für-Schritt-Anweisungen
    \item {\bf Abstraktion} -- Auf das Wesentliche fokussieren
    \item {\bf Musterkennung} -- Bekannte Lösungen wiederverwenden
\stopitemize

\subsection{Git Cheatsheet}

\starttabulate[|l|p(8cm)|]
\HL
\NC {\bf Befehl} \NC {\bf Beschreibung} \NC\NR
\HL
\NC \code{git init} \NC Neues Repository erstellen \NC\NR
\NC \code{git status} \NC Änderungen anzeigen \NC\NR
\NC \code{git add .} \NC Alle Änderungen stagen \NC\NR
\NC \code{git commit -m "..."} \NC Commit mit Nachricht \NC\NR
\NC \code{git log --oneline} \NC Kompakte Historie \NC\NR
\NC \code{git diff} \NC Änderungen anzeigen \NC\NR
\HL
\stoptabulate

\subsection{Prompt Engineering Cheatsheet}

{\bf Template:}
\startbashcode
Erstelle [Was] in [Sprache/Framework].

Kontext: [Warum/Wofür]

Anforderungen:
- [Anforderung 1]
- [Anforderung 2]

Beispiel:
Eingabe: [Beispiel]
Ausgabe: [Beispiel]

Einschränkungen:
- [Was zu beachten ist]
\stopbashcode

\blank[big]

\starttippbox
{\bf Erfolgskriterien für Modul 1}

Sie haben das Modul erfolgreich abgeschlossen, wenn Sie:
\startitemize[packed]
    \item Eine funktionierende Entwicklungsumgebung haben
    \item Einfache Programme mit KI-Unterstützung erstellen können
    \item Den Entwicklungsworkflow mit Git verstehen
    \item Effektive Prompts für Code-Generierung formulieren können
    \item KI-generierten Code auf Qualität überprüfen können
\stopitemize
\stoptippbox

\stopcomponent

